\section{Experimental Results}

To prove the validity of our response-time analysis, we went on running some
experiments on the software application executing on a real development board
(Olimex STM32\-H405). In particular, the goal was to validate the scheduling
feasibility of the system using the computed workloads from the RTA, and then
obtain the values that make the system fail to meet deadlines
(\textit{sensitivity analysis}). The latter experiments were only performed on
the Rust system. 

\subsubsection{Holistic Analysis}

From the holistic analysis, we obtained the workloads reported in table
\ref{tab:exp-holistic-baseline} that maximize the system utilization. 
\begin{table}
  \centering
  \resizebox{\linewidth}{!}{%
    \begin{tabular}{lcc}
      \toprule
      \textbf{Task} & \textbf{Workload (ops)} & \textbf{Scaling} \\
      \midrule
      \texttt{Regular\_Producer}       & 15120 & 756 $\times$ 20 \\
      \texttt{On\_Call\_Producer}      & 10008 & 278 $\times$ 20 $\times$ 1.8 \\
      \texttt{Activation\_Log\_Reader} & 6755  & 139 $\times$ 20 $\times$ 1.8 $\times$ 1.35 \\
      \bottomrule
    \end{tabular}
  }
  \caption{Holistic baseline workloads (board experiments)}
  \label{tab:exp-holistic-baseline}
\end{table}
With this computed workloads both systems managed to meet all the deadlines
after execution of 5 \textit{hyperperiods}. 

Then, we started experimenting with varying workloads by focusing on one task at
a time keeping the others fixed, we obtained the results in table
\ref{tab:exp-holistic-sensitivity}. Workloads values were incremented at steps
of 100. 

\begin{table}
  \centering
  \resizebox{\linewidth}{!}{%
    \begin{tabular}{lccc l}
      \toprule
      \textbf{Scenario} & \textbf{RP (ops)} & \textbf{OCP (ops)} & \textbf{ALR (ops)} & \textbf{Observed} \\
      \midrule
      Baseline          & 15120 & 10008 & 6755  & No deadline misses \\
      \midrule
      RP increase       & 15200 & 10008 & 6755  & ALR misses $\sim$1/4 (when released with OCP) \\
      RP increase       & 16000 & 10008 & 6755  & RP always misses; ALR misses as before \\
      \midrule
      OCP increase      & 15120 & 10100 & 6755  & ALR misses $\sim$1/4 (when released with OCP) \\
      OCP increase      & 15120 & 16900 & 6755  & OCP always misses; ALR misses as before \\
      \midrule
      ALR increase      & 15120 & 10008 & 6900  & ALR misses $\sim$1/4 \\
      ALR increase      & 15120 & 10008 & 16900 & ALR always misses; no impact on others \\
      \bottomrule
    \end{tabular}
  }
  \caption{Holistic sensitivity results (board experiments)}
  \label{tab:exp-holistic-sensitivity}
\end{table}

Results matched closely the static analysis performed, since even small
increments from the original baseline incurred in deadline misses. These results
also highlight a key property of FPS: under overload conditions the tasks gets
affected in priority order (from lower to highest). 

\subsection{Offset-based Analysis}

Using offset-based analysis we obtained the workloads in table
\ref{tab:exp-offset-baseline} that maximize the system utilization. 

\begin{table}
  \centering
  \resizebox{\linewidth}{!}{%
    \begin{tabular}{lcc}
      \toprule
      \textbf{Task} & \textbf{Workload (ops)} & \textbf{Notes} \\
      \midrule
      \texttt{Regular\_Producer}       & 15120 & Baseline from analysis \\
      \texttt{On\_Call\_Producer}      & 15973 & Baseline from analysis \\
      \texttt{Activation\_Log\_Reader} & 1     & Baseline from analysis \\
      \bottomrule
    \end{tabular}
  }
  \caption{Offset-based baseline workloads (board experiments)}
  \label{tab:exp-offset-baseline}
\end{table}

No deadline misses were observed after 10 \textit{hyperperiods}.

The sensitivity analysis produced the values in table \ref{tab:exp-offset-sensitivity}.

\begin{table}
  \centering
  \resizebox{\linewidth}{!}{%
    \begin{tabular}{lccc l}
      \toprule
      \textbf{Scenario} & \textbf{RP (ops)} & \textbf{OCP (ops)} & \textbf{ALR (ops)} & \textbf{Observed} \\
      \midrule
      Baseline          & 15120 & 1 & 15973     & No deadline misses \\
      \midrule
      RP increase       & 16000 & 1 & 15973     & RP and ALR always miss \\
      \midrule
      OCP increase      & 15120 & 900   & 15973     & ALR misses $\sim$1/4 (when released with OCP) \\
      OCP increase      & 15120 & 16900 & 15973     & OCP always misses; ALR misses as before \\
      \midrule
      ALR increase      & 15120 & 15973 & 16900 & ALR always misses; no impact on others \\
      \bottomrule
    \end{tabular}
  }
  \caption{Offset-based sensitivity results (board experiments)}
  \label{tab:exp-offset-sensitivity}
\end{table}

This time, results were slightly more distant from the ones obtained from static
analysis. This could be due to an overestimation of the WCET of system and
operation overhead. 

Nonetheless, results still follow a consistent behavior when overload is reached.


\section{Conclusions}
In this technical report, we presented the implementation of a real-time system
in Rust, and we compared it with an equivalent implementation in Ada. The
analysis, conducted with the \textit{MAST} tool, showed that both
implementations are able to achieve similar results in terms of system
utilization. To support the analysis, both systems were also tested on an
\textit{Olimex STM32\-H405} \cite{olimexSTM32H405} using the identified
workloads, and no missed deadlines were observed. Then we tested the validity of
our analysis, in terms of "pessimism" (sensitivity), obtaining consistent
results. \\ 
Rust, and in particular the \textit{RTIC} framework, proved to be a
valid alternative for the development of real-time systems. \\
Regarding the \textit{RTIC} framework, by being a declarative system model for
single-core applications, it enforces a well-defined structure,
requiring a clear understanding and strict control over the management of both
local and shared resources. This strong structure can support developers in
maintaining a well-organized and systematic design of the system. However, in
some situations it may also limit design flexibility in relation to \textit{Ada
Ravenscar Profile}, especially when dealing with shared resources. In
particular, we were unable to separate the responsibility of the
\texttt{on\_call\_producer} activation, as it was not possible to implement it
except by embedding it directly within the \texttt{deposit} call on the
\texttt{Request\_Buffer}. This limitation stems from the fact that any access to
a shared resource must be explicitly declared in the task’s shared resources
and, consequently, it is not possible for a task to provide an interface through
which other tasks may interact with it. The management of shared resources
always requires acquiring a lock when accessing them, which can introduce
unnecessary blocking times even in cases where the relationships between tasks
exclude the possibility of simultaneous access. Once again, this enforced
structure represents the trade-off between flexibility and the strong safety
guarantees provided by \textit{Rust}. \\
Unlike the \textit{Ada Ravensacar Profile}, on the subject of deadline miss
detections, the framework reveals a significant limitation by providing no
built-in support for handling such events, thus it leaves the responsibility for
handling timing violations entirely to the developers. \\
Our implementation are available at \\
\url{https://github.com/tommasoprandin/rtks} (\textit{Rust} system) and 
\url{https://github.com/tommasoprandin/gee} (\textit{Ada} system). Regarding
the identification of certain \textit{RTIC} system overheads, the results can be
found at \url{https://github.com/GianlucaBresolin/Profiling-RTIC}. \\  
Finally, we consider the \textit{RTIC} framework to be a valid, flexible and
performant option for the development of real-time systems, with a generated
code that, as shown in this report, infers no substantial overhead compared to
an equivalent implementation in \textit{Ada Ravenscar Profile}. 
